\documentclass[11pt]{article}% Shared preamble for the rigor documents (Lamport-style structured proofs).  \input this file after \documentclass.
\usepackage[T1]{fontenc}\usepackage{lmodern}\usepackage{microtype}
\usepackage[margin=1in]{geometry}
\usepackage{amsmath,amssymb,amsthm,mathtools}
\usepackage{xcolor}\usepackage{enumitem}\usepackage{booktabs}\usepackage{graphicx}\usepackage{hyperref}
\usepackage[most]{tcolorbox}
\definecolor{navy}{HTML}{17365D}\definecolor{teal}{HTML}{117A8B}\definecolor{warn}{HTML}{A33A2B}\definecolor{okgreen}{HTML}{2E7D4F}
\hypersetup{colorlinks=true,linkcolor=navy,citecolor=teal,urlcolor=teal}
\theoremstyle{plain}
\newtheorem{theorem}{Theorem}[section]\newtheorem{lemma}[theorem]{Lemma}\newtheorem{proposition}[theorem]{Proposition}\newtheorem{corollary}[theorem]{Corollary}
\theoremstyle{definition}\newtheorem{definition}[theorem]{Definition}\newtheorem{remark}[theorem]{Remark}\newtheorem{claim}[theorem]{Claim}\newtheorem{issue}[theorem]{Issue}
% ---------- Lamport structured proofs ----------
% \lstep{level}{number}{assertion}   -> indented "<level>number. assertion"
% \lproof                             -> "PROOF:" line (start of the sub-proof of the preceding step)
% \lby{justification}                 -> "BY ..." leaf justification for the preceding step
% \lqed{level}{number}                -> QED step of a sub-proof
% keywords: \lassume \lprove \lsuffices \lcase \ldefine \lpick \lnew
\newlength{\lindent}\setlength{\lindent}{1.7em}
\newcommand{\lstep}[3]{\par\smallskip\noindent\hspace*{\dimexpr\numexpr#1-1\relax\lindent\relax}{\color{navy}$\langle#1\rangle#2$.}\ #3}
\newcommand{\lproof}{\par\noindent\hspace*{\lindent}{\scshape Proof:}}
\newcommand{\lby}[1]{\par\noindent\hspace*{\lindent}{\scshape By}\ #1.}
\newcommand{\lqed}[2]{\par\smallskip\noindent\hspace*{\dimexpr\numexpr#1-1\relax\lindent\relax}{\color{navy}$\langle#1\rangle#2$.}\ {\scshape Q.E.D.}}
\newcommand{\lassume}{{\scshape Assume}\ }\newcommand{\lprove}{{\scshape Prove}\ }\newcommand{\lsuffices}{{\scshape Suffices}\ }
\newcommand{\lcase}{{\scshape Case}\ }\newcommand{\ldefine}{{\scshape Define}\ }\newcommand{\lpick}{{\scshape Pick}\ }\newcommand{\lnew}{{\scshape New}\ }
% ---------- verdict boxes ----------
\newtcolorbox{verdictok}{colback=okgreen!8,colframe=okgreen,title={\bfseries Verdict: verified},fonttitle=\small}
\newtcolorbox{verdictgap}{colback=orange!10,colframe=orange!80!black,title={\bfseries Verdict: gap / caveat},fonttitle=\small}
\newtcolorbox{verdictbreak}{colback=warn!10,colframe=warn,title={\bfseries Verdict: proof-breaking issue},fonttitle=\small}
\newtcolorbox{citedbox}[1]{colback=navy!5,colframe=navy!60,title={\bfseries Cited result: #1},fonttitle=\small,breakable}
\setlength{\parindent}{0pt}\setlength{\parskip}{4pt}\allowdisplaybreaks
\newcommand{\cA}{\mathcal A}\newcommand{\cF}{\mathcal F}\newcommand{\cM}{\mathfrak M}\newcommand{\cN}{\mathfrak N}

% extra calligraphic shortcuts (\providecommand: no clash if a document defines its own)
\providecommand{\cB}{\mathcal B}\providecommand{\cC}{\mathcal C}\providecommand{\cD}{\mathcal D}\providecommand{\cE}{\mathcal E}\providecommand{\cG}{\mathcal G}\providecommand{\cH}{\mathcal H}\providecommand{\cI}{\mathcal I}\providecommand{\cJ}{\mathcal J}\providecommand{\cK}{\mathcal K}\providecommand{\cL}{\mathcal L}\providecommand{\cO}{\mathcal O}\providecommand{\cP}{\mathcal P}\providecommand{\cQ}{\mathcal Q}\providecommand{\cR}{\mathcal R}\providecommand{\cS}{\mathcal S}\providecommand{\cT}{\mathcal T}\providecommand{\cU}{\mathcal U}\providecommand{\cV}{\mathcal V}\providecommand{\cW}{\mathcal W}\providecommand{\cX}{\mathcal X}\providecommand{\cY}{\mathcal Y}\providecommand{\cZ}{\mathcal Z}
\newcommand{\Fid}{\operatorname{F}}\newcommand{\Tr}{\operatorname{Tr}}\newcommand{\eps}{\varepsilon}\newcommand{\dd}{\,\mathrm d}

\newcommand{\hh}{\mathfrak h}
\newcommand{\Herm}{\operatorname{Herm}}
\newcommand{\SLD}{\mathcal S_Q}
\newcommand{\Aop}{\mathbb A}
\newcommand{\Stwo}{\mathfrak S_2}
\geometry{top=0.9in,bottom=0.9in,left=0.95in,right=0.95in}
\AtBeginDocument{\setlength{\parskip}{2.5pt}}
\newcommand{\vA}{\textcolor{okgreen}{\textsc{applied}}}
\newcommand{\vP}{\textcolor{orange!80!black}{\textsc{applied, residue}}}
\newcommand{\vF}{\textcolor{warn}{\textsc{not applied}}}
\title{Third (closing) referee report: \texttt{rigor/exact\_optimum\_tangent\_problem.tex}\\
\large (Phase 5, closing review: verification of the repairs S1--S6 and of the
statement and use of Hypothesis~H3)}
\author{Referee REF-P5-4 (Opus)}\date{2026-09-15}
\begin{document}\maketitle
\begin{abstract}
Closing review of F1's \texttt{exact\_optimum\_tangent\_problem.tex} (26\,pp, 1605 lines,
2026-09-15) after the repairs S1--S6 required by
\texttt{referee2\_exact\_optimum\_tangent\_problem.tex} \S\,``Required repairs''.  \textbf{All
six repairs are applied}, four without residue.  H3 is declared as Def.~4.12
(\texttt{hyp:H3}), Rem.~4.13 names the two excluded families, Prop.~4.10(3) is restricted to
vector-field generators, every ``(e) holds'' is hedged to ``on the model cone'', and
$\kappa_{\rm opt}=1-\theta=0.616\pm0.003$ is conditional on H1, H2 \emph{and} H3 in
Rem.~4.11, \S7 and the \texttt{verdictgap} box.  The
use of H3 is sound where it matters (\S\ref{sec:h3}).  Three residual defects remain,
none of which touches the conclusion: (D1) the ``equivalently'' clause of Def.~4.12 is
\emph{not} equivalent to (33) and is trivially satisfied by the second family
Rem.~4.13 excludes; (D2) Thm.~3.11(3) is still \emph{stated} with ``$\mathcal G$ \dots\ an
isometry generator'' although its repaired proof step $\langle1\rangle4$ declines to prove
exactly that; (D3) the ambient space of (33) is unspecified, and read inside
$\mathfrak F$ it forces $\alpha=0$, which makes the advertised collapse of (e) to the boxed
(31) degenerate.  Verdict: \textsc{accepted}; D1 and D2 are one-line
corrections, D3 is a clarification.
\end{abstract}

\section{Scope and method}\label{sec:scope}
Read in full against \texttt{referee2\_exact\_optimum\_tangent\_problem.tex} \S``Required
repairs'' (S1--S6) and the author's account in \texttt{PHASE5\_F1\_SUMMARY.md} rev.~4.
Locators taken from the current \texttt{exact\_optimum\_tangent\_problem.aux}: the
numbering is unchanged through Rem.~4.11, the new material is Def.~4.12 (\texttt{hyp:H3},
p.~20), eq.~(33) and Rem.~4.13 (\texttt{rem:H3}, p.~20); the document
compiles at 26\,pp.  Inside Prop.~4.10 the manual label (2$'$) does not advance the
counter, so the rendered items are 1, (2$'$), 2, 3 (two model directions),
4 ($\mathcal E(\zeta_\infty)>0$), 5 (model cone), and \S5--\S7 cite them correctly (checked
against the PDF).  I edited neither the document nor the earlier referee reports.

\section{The six repairs}\label{sec:s16}
{\footnotesize\begin{tabular}{@{}llp{9.9cm}@{}}
\toprule
 & Status & Where, and what is left\\
\midrule
S1 & \vP & H3 is Def.~4.12 in the form ``for every admissible $\zeta$ with
 $\Aop(\zeta,0,0)\in\mathfrak F$ there are $\alpha,\beta\ge0$, $\eta\in\overline L$ with
 (33)''; Rem.~4.13 names the two excluded families (multiplicity-$>1$ shifts via
 Cooper--Wold, correctly noting that the dilation $f\mapsto\sqrt2f(2y)$ is \emph{not} the
 sharp example; unbounded skew-adjoint $\zeta=\mathrm im(y)$, for which (28) is unavailable
 because the cyclicity rearrangement needs $\zeta$ bounded).  Prop.~4.10(3) now considers
 ``the generators of \emph{vector-field} type $\zeta_v=-D_v$'' and calls the result the
 ``\emph{vector-field subcone} --- \emph{not}, as far as is proved here, the whole
 admissible cone''.  All eleven
 occurrences of ``(e) holds'' are hedged to the model cone (Prop.~4.10(5), Rem.~4.11,
 Rem.~4.13, \S5, \S7 Proved~5 / Numerical~3 / Conditional~3--4, \texttt{verdictgap});
 $\kappa_{\rm opt}=1-\theta$ carries ``conditional on H1, H2 \emph{and} H3'' in all four
 places it is asserted.  The abstract declares H3 and makes no unhedged
 $\kappa_{\rm opt}$ claim.  Residue: D1, D3.\\
S2 & \vA & Rem.~4.11(ii) gives the decomposition (32), states that the
 total is dominated by the upwind-viscosity middle term and grows like $h^{-0.97}$ (ratios
 $1.960,1.963,1.965$) whereas the corner grows like $h^{-0.60}$ (ratios
 $1.562,1.533,1.508$), calls the total \emph{scheme-dependent}, names the corner row
 $\frac12M_*(0,0)/h=+10.3,+16.1,+24.7,+37.2$ the continuum quantity and the row whose
 margin is to be quoted.  ``Driven by the corner term'' survives only as a negated
 quotation; \S7 Numerical~3 repeats the corrected reading (``must not be quoted as the
 margin''), and ``$+\infty$ is an allowed value'' is gone.\\
S3 & \vA & (27) has values in $[-\infty,+\infty]$, its second case flagged as a
 \emph{convention, not an extension of the form} $M_*$, with the statement that
 such $\zeta$ carry no feasible first-order datum (Def.~3.1(3)) so that (e) is
 \emph{vacuous} there, and that reading $+\infty$ as a margin is legitimate only after (a).
 The boxed (31) is split: the inequality $\mathcal E(\zeta_0)\ge0$, then
 $\mathcal E(\zeta_0)=\frac12\Tr((-2\sigma_{\zeta_0})M_*)$ \emph{provided}
 $-2\sigma_{\zeta_0}$ lies in the form domain of $M_*$ --- a proviso declared to fail
 ($\delta_0\otimes\delta_0$ unbounded, $\delta_0\notin\hh_D$).  Def.~4.9 adds the two
 conventions: $\mathfrak D$ read as a linear subspace (equivalently, polarisation only for
 $\psi\pm\varphi\in\mathfrak D$).\\
S4 & \vP & Thm.~3.11(3) proof $\langle1\rangle4$ takes N4's second option: ``$\mathcal G$
 is read, by definition, as \emph{the anti-symmetric part of an admissible} $\zeta$ --- one
 may not conclude from `$\mathcal G$ and $\zeta$ agree on $C_c^\infty(D)$' that $\mathcal G$
 generates an isometry semigroup, since that needs $C_c^\infty(D)$ to be a core''.  The non
 sequitur is gone.  Residue: D2 (the \emph{statement} of (3) not adjusted to match).\\
S5 & \vA & \S7 Proved item~5 now opens ``the derivative condition along \emph{the
 non-isometric directions} $\Aop(0,N_1,Y_1)$ is (a) \dots, equivalently (c) \dots'' and
 only then gives global optimality as (a) \emph{together with} (e).  The refuted
 equivalence is no longer stated first, or at all.\\
S6 & \vA & (i) Thm.~4.4 $\langle1\rangle4$: ``on $\mathcal C_\zeta$ it is
 $\mathcal E(\zeta)$ itself'' (the spurious $\frac12$ gone).  (ii) Lem.~4.2
 $\langle1\rangle6$: $\Aop(\mathcal G,0,0)=[\mathcal G,Q]=-[Q,\mathcal G]$.
 (iii) Prop.~4.10(2$'$) quotes the four ladder values $-9.1\cdot10^{-6}$, $+1.9\cdot10^{-4}$,
 $-9.9\cdot10^{-3}$, $+7.4\cdot10^{-2}$ and the $O(1)$ outlier $-6.3$ (``not merely small
 noise'').  (iv) Rem.~6.4: $1-2/\pi$ is ``still $6.9\sigma$ low''.  (v) The
 $\theta$-problem (42) (formerly (36)) reads $\inf_{\mathcal G}$ over $\mathcal G$ bounded
 anti-self-adjoint supported in $\{y>0\}$, ``the class of Cor.~3.10''.\\
\bottomrule
\end{tabular}}

\section{Is H3 stated so that it is used correctly?}\label{sec:h3}
The use is Rem.~4.13: ``Under H3, $\mathcal E(\zeta)=\alpha\mathcal E(\zeta_0)
+\beta\mathcal E(\zeta_\infty)$ by Prop.~4.10(1), so condition (e) collapses to (31)''.

\paragraph{What it needs, and what is available.}  Apply $2b(\delta_*,\cdot)$ to
(33).  Two things are required, and both are present.  (i) \emph{Linearity in the
right slot of $b$}, not linearity of $\mathcal E$ on the cone of \emph{generators}: since
H3 is formulated at the level of $\Aop(\zeta,0,0)$ and not of $\zeta$, the identity follows
from bilinearity of the polarisation $b$ of the quadratic form $g_Q$ alone.  This is the
right way round: had H3 been stated as ``$\zeta=\alpha\zeta_0+\beta\zeta_\infty+\mathcal G$
modulo $\overline L$'' one would also have needed additivity of $\zeta\mapsto\Aop(\zeta,0,0)$
on sums of \emph{unbounded} generators with different domains, which is nowhere established.
(The citation of Prop.~4.10(1) is therefore imprecise --- bilinearity of $b$ is what is used
--- but harmless.)  (ii) \emph{The normal equation on the closure}, $b(\delta_*,\eta)=0$ for
$\eta\in\overline L$: available, and this is the point that could have failed.  $\overline L$
is defined at (19) as the closure of $L$ \emph{in the Hilbert space} $(\mathfrak F,b)$,
$\delta_*$ is the $b$-orthogonal residual onto that closed subspace, and
$|b(\delta_*,\eta)|\le g_Q(\delta_*)^{1/2}g_Q(\eta)^{1/2}$ (Cauchy--Schwarz, used already in
Def.~4.9) makes $b(\delta_*,\cdot)$ $g_Q$-continuous, so orthogonality extends from $L$ to
$\overline L$ automatically: no continuity hypothesis is smuggled in.

\paragraph{D1 (defect): the ``equivalently'' of Def.~4.12 is not an equivalence.}  Def.~4.12
adds ``equivalently, modulo $\overline L$ the boundary form $-2\sigma_\zeta$ is a positive
combination of the two model boundary forms of Prop.~4.10(3)''.  This is strictly weaker
than (33), and weaker exactly on the family Rem.~4.13(2) excludes: for
$\zeta=\mathrm im(y)$ one has $\sigma_\zeta=0$, so the boundary-form condition holds
trivially with $\alpha=\beta=0$, whereas (33) then demands
$\Aop(\zeta,0,0)\in\overline L$ --- precisely what Rem.~4.13(2) says ``may fail''
(Rem.~4.8(ii)).  The boundary-form reading would thus make H3 \emph{provable} for that family
without delivering $\mathcal E(\zeta)\ge0$ there, breaking H3~$\Rightarrow$~(e).  By the
absorption identity $\Aop(\zeta,0,0)=\Aop(\mathcal G_\zeta,-2\sigma_\zeta,0)$ the two
readings agree only under the proviso $\Aop(\mathcal G_\zeta,0,0)\in\overline L$, unavailable
for unbounded $\mathcal G_\zeta$.  \emph{Repair}: delete the clause or append that proviso.

\paragraph{D3 (clarification): the ambient space of (33), and a degenerate collapse.}
(33) contains $\Aop(\zeta_0,0,0)$, which the document's own numerics put \emph{outside}
$\mathfrak F$ (Rem.~4.11(ii), Prop.~4.10(5)), while its left-hand side, $\eta\in\overline L$
and $\Aop(\zeta_\infty,0,0)=\lambda\dot\delta$ are all in $\mathfrak F$.  Read inside it --- the
document's global convention, ``\emph{All} statements below are inside $\mathfrak F$'' (after
(4)) --- (33) forces $\alpha=0$; read as an identity of sesquilinear forms on $C_c^\infty(D)$,
as (14) is, it admits $\alpha>0$, but then Rem.~4.13's displayed identity has a finite
left-hand side and $+\infty$ on the right, so again $\alpha=0$.  Either way (e) on the H3
class follows from Prop.~4.10(4) ($\mathcal E(\zeta_\infty)>0$) alone, and the boxed (31) is
never invoked.  The \emph{sign} conclusion is unaffected --- under H3, (e) holds --- so \S7
and the \texttt{verdictgap} box need no change; what should change is the advertisement that
(e) ``collapses to'' the scalar test (31), vacuous under the document's own numerics.
\emph{Repair}: say in which space (33) is read, and note that with
$\Aop(\zeta_0,0,0)\notin\mathfrak F$ the hypothesis has content only through $\beta$.

\paragraph{D2 (defect): Thm.~3.11(3), statement versus repaired proof.}  The statement of
(3) still reads $\mathcal C\subseteq\{\Aop(\mathcal G,N_1',Y_1):\mathcal G$ anti-symmetric
\emph{and an isometry generator}$\}$, while the repaired $\langle1\rangle4$ says that one
may \emph{not} conclude that $\mathcal G=\zeta^*-\sigma$ generates an isometry semigroup, and
reads $\mathcal G$ on the right as the anti-symmetric part of an admissible $\zeta$.  So the
proof establishes the inclusion into the \emph{larger} set, not the displayed one.  Nothing
downstream breaks --- Thm.~4.4 $\langle1\rangle4$ and Rem.~4.8 use (3) only as an inclusion
that must not be enlarged to \emph{all} anti-symmetric $\mathcal G$, true under either
reading --- but the statement (and \S7 Proved item~4, which inherits the wording) should be
brought into line with $\langle1\rangle4$.

\emph{Minor}: (i) the \texttt{verdictgap} box calls F2 T5 ``positive and divergent at
$\zeta_0$'' without the S3 reading (vacuous, not a margin) that Rem.~4.11 and \S7
Numerical~3 give.  (ii) In S6(iii) one transcribed magnitude is wrong, inherited from N6(iii)
of the previous report: the outlying rungs of \texttt{f2\_t5\_endpoint.out} are
$-6.3146\cdot10^{1}$ and $-1.5721\cdot10^{6}$, so Prop.~4.10(2$'$)'s ``other rungs of the
scan reaching $-6.3$'' should read $-63$ (and may as well quote $-1.6\cdot10^{6}$) --- the
slip understates the author's own point.

\section{Verdict}\label{sec:verdict}
S1--S6 are \textbf{all applied}; S2, S3, S5 and S6 without residue.  The document is
\textsc{accepted}.  H3 is correctly placed: stating it at the level of $\Aop(\zeta,0,0)$ is
what makes its use legitimate, and both ingredients of that use --- bilinearity of $b$, and
the normal equation on $\overline L$, available because $\overline L$ is closed in
$(\mathfrak F,b)$ --- are present.  The residues are D1 (delete or qualify the
``equivalently'' clause of Def.~4.12), D2 (align the statement of Thm.~3.11(3) with its
repaired proof step), D3 (name the ambient space of (33)).  None changes an assertion of
\S2--\S4 or \S6, and none weakens the Phase-5 conclusion: numerically
$\kappa_{\rm opt}=1-\theta=0.616\pm0.003$, conditional on H1, H2, H3, with (a) and (e)
verified on meshes and only the moving-endpoint half of (e) proved.

\paragraph{Addendum (2026-09-15, after the repair of D1--D3; document now 27\,pp).}
All three residues are repaired and the minor slip with them.  \textbf{D1}: the
``equivalently \dots\ boundary form'' clause of Def.~4.12 is deleted, and the reason is now
stated there --- a boundary-form version would hold trivially for $\zeta=\mathrm im(y)$
($\sigma_\zeta=0$) while (33) demands $\Aop(\zeta,0,0)\in\overline L$, so it would not imply
(e).  \textbf{D3}: (33) is now read \emph{in $\mathfrak F$}, and the new Rem.~4.13
(\texttt{rem:H3space}) draws the consequence I asked for and draws it correctly:
$\Aop(\zeta_0,0,0)\notin\mathfrak F$ forces $\alpha=0$, so H3 says that every admissible
generator carrying a feasible datum is, modulo $\overline L$, a non-negative multiple of
$\zeta_\infty$; (e) then holds \emph{outright} under H3 through
$\mathcal E(\zeta_\infty)>0$, the junction inequality (31) is never invoked, and (31)
survives as the test for \emph{regularised} junction directions --- which is exactly what F2
T5 measures.  The same reading is carried into Prop.~4.10(5), \S5 item~5, \S7 (Proved~5,
Conditional H3) and the verdict box.  \textbf{D2}: Thm.~3.11(3) now describes the
right-hand side by \emph{how} $\mathcal G=\zeta^*-\sigma$ arises, ``not by the property of
being itself an isometry generator, which $\langle1\rangle4$ does not prove'', so statement
and proof agree.  Minor: Prop.~4.10(2$'$) now reads $-63$ and $-1.6\cdot10^{6}$.  Nothing
overclaims: ``(e) holds'' is conditional on H3 wherever it is unqualified, H3 is still
``proved for vector-field generators, open in general'', and
$\kappa_{\rm opt}=1-\theta=0.616\pm0.003$ remains conditional on H1, H2 and H3 in every
place it is asserted.  \emph{One leftover sentence}: Rem.~4.14 (\texttt{rem:H3}) still opens
``Under H3, $\mathcal E(\zeta)=\alpha\mathcal E(\zeta_0)+\beta\mathcal E(\zeta_\infty)$
\dots\ so condition (e) collapses to (31)'', which is the superseded formulation and
contradicts Rem.~4.13(1) (``(e) should not be advertised as `collapsing to' the scalar
test'') --- its closing sentence was updated but its opening one was not; it should read
``under H3, $\alpha=0$ and $\mathcal E(\zeta)=\beta\mathcal E(\zeta_\infty)\ge0$''.  With
that one clause the document is clean; the verdict of this report stands as
\textsc{accepted}.

\end{document}
